\section{What institutions and regional experts can explain}
\label{sec:expert-findings}
The Colombian institutional memo describes problems at particular points in service delivery, while the regional expert memo considers what differs across providers and countries. We use the two for implementation diagnoses, not for population estimates or as witnesses to the migrant accounts. The combined institutional report draws on those memos; a point repeated there does not count as another source. An official document serves a different purpose: it establishes a written rule or feature for the scheme it covers.

\subsection{Recognition, verification and keeping an account usable}

Colombian institutional interviews distinguish formal recognition of migration documents from verification in a provider's system. The memo raises inconsistent checks and difficulty preserving account access or financial records when documentation changes. It also describes problems that have little to do with legal eligibility: shared or damaged devices, changes in phone numbers, limited data, and trouble authenticating or navigating an app. Opening an account and maintaining usable access are different tasks. The accounts of rejected documents and borrowed access in Section~\ref{sec:migrant-findings} show what users report when a step fails; the institutional interviews help locate plausible points of failure. They do not establish which process failed in those particular cases or how often it fails.

Regional experts also point to variation in prior digital experience and in what an account or wallet allows a user to do. Some migrants arrive with experience using digital payments. Others may have an account but still need help at registration, or may face limits when seeking to transfer or withdraw funds. That evidence argues against treating all operational barriers as a skills deficit. Assistance at enrollment, a usable interface and continued access when a number or document changes address different problems. These are proposals to evaluate, not observed effects of training or redesign.

\subsection{A payment rule covers a step, not the entire journey}

The Colombian and regional accounts distinguish two jobs. Domestic interoperation connects participating providers; a transfer abroad may also require identity checks across jurisdictions, conversion, delivery and a way to resolve disputes. Standardizing the local step need not change who is responsible for later steps. The regional memo's concerns about limited active wallet use do not establish a contradiction with the high rate of reported domestic transactions in our survey: wallet use is not the same outcome, and the regional evidence has no comparable denominator.

The difference is concrete for complaints. For a domestic Bre-B transaction, Banco de la República directs users to the participating financial institution's complaint channel \citep{banrepfaq}. That identifies a point of entry for a dispute under the domestic scheme, not the party responsible for a foreign leg or an intermediary's actions. The Colombian and regional memos raise difficulties exercising remedies and coordinating across organizations. Neither the written route nor those diagnoses establish that someone in the study received a response or recovered money. The benchmark compares formal arrangements for precisely this purpose: to identify which design or institution could address a given step, not to rank realized access or protection across countries.
